\subsection*{Chapter 18: Improving GRPO}

\paragraph{Exercise 18.1} All three edits give 0: a typo makes the tag a different token
sequence, so \code{index} raises; reversed order fails the \code{<} test; a missing tag raises.
The string check catches the same three. The id-based check can be broken by tokenisation:
if the tokenizer does not treat \verb|<think>| as a special token (the Qwen3 \emph{base}
tokenizer splits it into three pieces), the text contains the tag but the single id never
occurs, so the id check returns 0 while the string check returns 1. Conversely, the model may
emit pieces that decode to the tag without the special id.

\paragraph{Exercise 18.2} In our runs (two seeds, 150 steps), every format variant drove
\code{format\_avg} to $\approx$\,0.99 within about 30 steps, because the SFT model already emitted
\verb|<think>| 20\% of the time and the tag is the easiest reward available. Unconditional
weight 1 reached a late-run \code{correct\_avg} of about 0.74 and greedy accuracy of 0.84--0.87,
below the baseline (0.83 and 0.89--0.96). Weight 0.1 behaved the same, since even a small
weight breaks ties between otherwise identical rollouts. The conditional variant recovered
the baseline's correctness (0.80--0.86, greedy 0.91--0.96). In \cref{fig:grpo2:formatfig},
under plain summation a wrong-but-tagged rollout ties with a right-but-untagged one, so early
gradient mainly teaches tagging. Conditioning makes correctness a prerequisite for any reward.

\paragraph{Exercise 18.3} With \code{inner\_epochs}$=1$ the clip fraction is exactly 0
(every ratio is 1). In our \code{clip\_reuse4} runs (4 full-batch passes, learning rate
$3\times10^{-4}$) the late-run clip fraction was only 2--3\%, yet the clipped run kept the
baseline's accuracy (0.88--0.95) while the unclipped one degraded to 0.50--0.62. A small
fraction of clipped tokens can carry most of the damage, because they are exactly the tokens
whose probability is running away. The clip fraction grows with the number of passes and with
the learning rate; where it crosses 0.2 depends on both, so measure it. Wall-clock time per
update falls roughly as $(t_{\text{sample}}+k\,t_{\text{update}})/k$, so the saving saturates
once $k\,t_{\text{update}}$ dominates. On real reasoning workloads, where sampling dominates,
$k=2$--4 is the usual sweet spot.

\paragraph{Exercise 18.4} $\E_p[r]=\sum_yp(y)\frac{q(y)}{p(y)}=1$, so
$\E_p[k_3]=1-1-\E_p[\log r]=\E_p[\log\frac{p}{q}]=\KL(p\|q)$. Pointwise, $f(r)=r-1-\log r$
has $f(1)=0$, $f'(r)=1-1/r$ and $f''>0$, so $f\ge0$. Numerically, all three sample means
converge to the true KL ($k_2$ to a slightly different value, being biased). $k_1$ has
the largest variance and negative samples; $k_3$ typically has the smallest variance when $p$
and $q$ are close.

\paragraph{Exercise 18.5} With zero advantages, the sequence surrogate's gradient is
$\beta\nabla\sum_t\log\pitheta(y_t)$ on the fixed sampled tokens. Descending it lowers the
probability of exactly those tokens, so probability flows to everything else, the
sampled-token probability falls and entropy rises monotonically. The per-token $k_3$ term has
gradient $\beta(1-r_t)\nabla\log\pitheta(y_t)$ with $r_t=\piref/\pitheta$. It is zero when the
policy equals the reference and \emph{changes sign} when the policy undershoots
($r_t>1$ raises the token's probability again). It therefore pulls toward $\piref$
rather than away from the samples, and entropy returns to the reference's level.

\paragraph{Exercise 18.6} Per-token gradient weights for (short $+1$, long $-1$):
\begin{center}\small
\begin{tabular}{@{}lccc@{}}
\toprule
rule & per token (short) & per token (long) & total long\\
\midrule
\code{seq\_mean\_token\_mean} & $1/20=0.05$ & $1/2000=0.0005$ & $0.5$\\
\code{token\_mean} & $1/1010$ & $1/1010$ & $0.99$\\
\code{seq\_mean\_token\_sum\_norm} ($L=1000$) & $1/2000$ & $1/2000$ & $0.5$\\
\bottomrule
\end{tabular}
\end{center}
Per-sequence averaging lets each token of the long wrong answer off 100 times more lightly
than a token of the short right answer, which over many updates favours long wrong answers.
Token-level and constant normalisation treat tokens equally.

\paragraph{Exercise 18.7} With std scaling, $\sum_i|A_i|=\frac{2Gp(1-p)}{s}
=2\sqrt{G(G-1)}\sqrt{p(1-p)}$. Without it, $2Gp(1-p)$. The ratio (with/without) is $1/s$,
which is largest when $p(1-p)$ is small, i.e.\ for nearly-always-solved and nearly-never-solved
prompts. For $G=8$: at $p=\frac18$ the scaled weight is 4.95 versus 1.75 unscaled; at
$p=\frac12$ it is 7.48 versus 4. Std scaling up-weights edge cases relative to medium-difficulty
prompts. Dr.~GRPO calls this a difficulty bias. One could argue the reverse, that rare
successes on hard prompts are the most informative, which is why DeepSeek-V3.2 keeps it.

\paragraph{Exercise 18.8} Full \code{token\_entropy} on a $(4096,\num{151936})$ tensor creates
float32 copies of the logits, log-probs and their product: about three times 2.5\,GB. The chunked
version holds $\approx3\times C\times|\mathcal V|\times4$\,B, i.e.\ $\approx$\,0.47\,GB at
$C=256$ and $\approx$\,1.9\,GB at $C=1024$, with run time almost unchanged above
$C\approx128$ (each chunk is still a large, efficient kernel). Below that, launch overhead dominates.

\paragraph{Exercise 18.9} The policy with AdamW needs $32\times10^9\times16\,\text{B}/16=32$\,GB
per GPU when FSDP-sharded. (a) A full bf16 reference adds 64\,GB per GPU, 96\,GB in total,
which does not fit on 80\,GB GPUs. (b) Sharding it like the policy adds 4\,GB per GPU, plus an
all-gather of $\approx64$\,GB of reference weights per forward pass (each GPU receives
$\frac{15}{16}\cdot64$\,GB per step): $\approx0.15$\,s at 400\,GB/s NVLink, over a second across
50\,GB/s InfiniBand, overlappable layer by layer. (c) The rollout workers compute reference
log-probs. That costs trainer memory nothing and ships only 4 bytes per token (4\,MB per million
tokens), but the rollout fleet must hold an extra 64\,GB copy (or reuse the base weights if the
policy is a LoRA adapter). Option (c) is common in disaggregated systems, and (b) in colocated ones.

\paragraph{Exercise 18.10} Casting to bf16 changes per-token log-probs by roughly
$10^{-3}$--$10^{-2}$ nats, with heavy tails on low-probability tokens. On the toy task,
training is stable either way, because sequences are short and the mismatch is small. The
truncated weights mostly sit within $[0.99,1.01]$, with a few capped outliers. The effect
becomes important on long sequences (thousands of tokens, where small per-token differences
compound) and with aggressive sampler optimisations such as FP8 KV caches or speculative decoding.
